\documentclass{article}
\usepackage{graphicx}
\usepackage[russian]{babel}

\title{Мини-конспект по теме: Теорема Пифагора}
\author{Я Сама}
\date{Октябрь 2026}

\begin{document}

\maketitle{}
\tableofcontents{}

\section{Введение}
Теорема Пифагора — одна из важнейших теорем евклидовой геометрии. Она находит
применение в самых разных областях:
\begin{itemize}
    \item геометрия и тригонометрия
    \item физика
    \item инженерные расчёты
    \item компьютерная графика
\end{itemize}
\section{Формулировка теоремы}
\textbf{Слова:} В прямоугольном треугольнике квадрат гипотенузы равен сумме квадратов
катетов.

\begin{equation}
c^2=a^2+b^2
    \label{sec:1}
\end{equation}
\begin{center}
 Как видно из формулы \ref{sec:1}, знание двух сторон позволяет найти третью.   
\end{center}
\section{Доказательство (набросок)}
\begin{flushleft}
    Одно из доказательств основывается на площади квадрата, составленного из четырёх одинаковых прямоугольных треугольников и малого квадрата в центре. Раскладывая площадь двумя способами, получаем $c^2=a^2+b^2$.
\end{flushleft}


\end{document}
